\def\Tamb{85}       % Temperatura ambiente: 85
\def\Pv{4}          % Presión de vapor actual: 40

\def\Tas{60.0}  
\def\Pas{20}

% --- water_saturation_adiabatic.tex ---
\begin{tikzpicture}
\begin{axis}[
    width=11cm, height=8.5cm,
    ymode=log, 
    log ticks with fixed point,
    title={\textbf{Proceso de Saturación Adiabática ($H_2O$)}},
    xlabel={$\temperature[]{}$ [$^\circ$C]},
    ylabel={$\pressure[]{\vapor}$ [kPa] (Log)},
    xmin=\Tmin, xmax=\Tmax, 
    ymin=\Pmin, ymax=\Pmax,
    axis lines=left,
    % Ticks específicos para resaltar Tas y Tamb
    xtick={0, \Tas, \Tamb}, 
    xticklabels={0, $\temperature[]{sat}$, $\temperature[]{amb}$},
    ytick={0.1, \Ptp, 10, \Pv, \Pas, 101.3}, 
    yticklabels={0.1, \Ptp, 10, $\pressure[]{\vapor, 1}$, $\pressure[]{\text{sat}}(\temperature[]{sat})$, 101.3},
    grid=both,
    minor grid style={gray!5},
    major grid style={gray!15},
    clip=true
]

    % --- CURVAS DE COEXISTENCIA ---
    % Sublimación
    \addplot[teal, ultra thick, domain=\Tmin:\Ttp, samples=100] {0.611*exp(22.5*(1 - 273.16/(x+273.15)))};
    
    % Fusión
    \draw[blue, thick] (axis cs:\Ttp,\Ptp) -- (axis cs:-1.5,\Pmax); 
    
    % Vaporización (Saturación)
    \addplot[blue, ultra thick, domain=\Ttp:\Tmax, samples=100] {0.611*exp(17.27*x/(x+237.3))}
        node[pos=0.3, sloped, above, black, font=\tiny] {Curva de saturación, $\pSat$};

    % --- ETIQUETAS DE ESTADO ---
    \node[blue!20!black, font=\footnotesize] at (axis cs:-15, 50) {\textbf{SÓLIDO}};
    \node[blue!70!black, font=\footnotesize] at (axis cs:35, 80) {\textbf{LÍQUIDO}};
    \node[orange!70!black, font=\footnotesize] at (axis cs:45, 0.5) {\textbf{VAPOR}};

    % --- PROCESO DE SATURACIÓN ADIABÁTICA ---
    % Trayectoria del estado 1 (ambiente) al estado 2 (saturado adiabáticamente)
    \draw[red, ultra thick, -{Stealth[scale=1.2]}] (axis cs:\Tamb,\Pv) -- (axis cs:\Tas,\Pas)
        node[midway, above, sloped, font=\tiny, text=red!50!black] {Sat. Adiabática};
    
    % Proyecciones de los puntos de interés
    % \draw[gray, dashed] (axis cs:\Tas,\Pas) -- (axis cs:\Tas,\Pmin);
    % \draw[gray, dashed] (axis cs:\Tamb,\Pv) -- (axis cs:\Tamb,\Pmin);
    % \draw[gray, dashed] (axis cs:\Tas,\Pas) -- (axis cs:\Tmin,\Pas);
    % \draw[gray, dotted] (axis cs:\Tamb,\Pv) -- (axis cs:\Tmin,\Pv);

    % Punto 1: Aire de entrada (Estado inicial)
    \node[circle, fill=black, inner sep=1.5pt, label=right:{\scriptsize 1}] at (axis cs:\Tamb,\Pv) {};
    
    % Punto 2: Aire saturado a la salida (Estado final)
    \node[circle, fill=red, inner sep=1.5pt, label=above left:{\scriptsize 2}] at (axis cs:\Tas,\Pas) {};

    % Punto Triple
    \filldraw (axis cs:\Ttp,\Ptp) circle (1.5pt) node[right, font=\tiny, xshift=2pt] {\triplePoint};

\end{axis}
\end{tikzpicture}
